\section{Introduction}
\label{sec:intro}

Sustaining human exploration beyond Low Earth Orbit (LEO) requires the establishment of self-sustaining Bioregenerative Life Support Systems (BLSS) \citep{fu2016how}. Central to BLSS architecture are higher plants, which perform vital functions including carbon dioxide sequestration, oxygen generation, wastewater purification through transpiration, and dietary supplementation with fresh nutrients \citep{paul2017genetic,dixit2021persistence}. For decades, the model crucifer \textit{Arabidopsis thaliana} has served as the foundational benchmark for elucidating how plants perceive and adapt to extraterrestrial environments, notably altered gravity vectors, chronic low-dose cosmic radiation, and fluid-behavior anomalies \citep{choi2019variation,kruse2020spaceflight}.

Extensive transcriptomic and proteomic profiling archived in the NASA Open Science Data Repository (OSDR) has illuminated pervasive reprogramming of plant stress responses, hormone signaling, cell-wall remodeling, and epigenetic modifications during orbital flight \citep{barker2023meta,zhou2019epigenomics,olanrewaju2023integrative}. However, terrestrial plant biology has firmly established that plants do not function as isolated genetic entities; rather, they exist as ``holobionts'' intrinsically co-adapted with diverse microbial consortia comprising endophytic, phyllospheric, and rhizospheric bacteria and fungi \citep{dezelicourt2013rhizosphere,pieterse2014induced}. Beneficial endophytes confer pathogen resistance through induced systemic resistance (ISR), synthesize phytohormones such as indole-3-acetic acid (IAA), and alleviate abiotic stressors including drought and oxidative challenge. Conversely, alterations in host immunity in microgravity may render spaceflight crops vulnerable to opportunistic colonization by phytopathogens or human bacterial pathogens \citep{totsline2023microgravity,totsline2024simulated}.

Despite its paramount importance for space agriculture, characterization of the plant microbiome during spaceflight has faced significant experimental constraints. Launch mass limits, astronaut crew time limitations, and strict containment protocols often restrict dedicated metagenomic sampling. Consequently, prior studies of the International Space Station (ISS) microbiome have primarily focused on environmental air filters, surface swabs, and astronaut skin or gut microbiomes, leaving seedling-associated endophytic communities under-sampled.

Fortunately, modern high-throughput RNA sequencing (RNA-seq) performed on total plant tissue inherently captures a rich fraction of non-host transcripts \citep{park2019systematic}. When plant total RNA is extracted and sequenced without stringent poly(A) selection or when microbial transcripts share sequence features that evade depletion, hundreds of thousands of sequencing reads originate from bacteria, archaea, and fungi inhabiting the host tissues. In terrestrial genomics, such reads are traditionally categorized as ``unmapped dark matter'' and discarded during downstream differential expression analyses \citep{park2019systematic,ewels2020nf}.

In this study, we capitalize on this bioinformatic reservoir by conducting a systematic meta-transcriptomic recovery and microbial profiling of archival \textit{Arabidopsis thaliana} spaceflight datasets deposited in NASA OSDR \citep{barker2023meta}. By comparing multiple flight experiments conducted in diverse hardware—including the Biological Research in Canisters (BRIC-19, BRIC-20, BRIC-PDF) and the Advanced Plant Experiments (APEX-03-2) hardware—across four distinct natural accessions (Col-0, WS-2, Ler-0, and Cvi-0), we address fundamental questions regarding spaceflight microbiome dynamics:
\begin{enumerate}
    \item Do surface-sanitized \textit{Arabidopsis} seedlings grown in closed spaceflight hardware harbor persistent, detectable microbial communities?
    \item How are these communities partitioned between intrinsic seed-borne endophytes and extrinsic human or hardware commensals?
    \item Does spaceflight microgravity or hardware containment elicit consistent taxonomic shifts relative to matched ground controls?
    \item How do host developmental stage, tissue type (whole seedling vs.\ dissected root tip), and ecotype dictate microbial composition in orbit?
\end{enumerate}
